% Generated draft: requires mathematical and source review.
\documentclass[11pt]{article}
\usepackage[margin=25mm]{geometry}
\usepackage{fontspec}
\setmainfont{Latin Modern Roman}
\usepackage{amsmath,amssymb,mathtools,unicode-math}
\setmathfont{Latin Modern Math}
\usepackage{enumitem,needspace,xurl,hyperref}
\setlength{\parindent}{0pt}
\setlength{\parskip}{5pt}
\setlength{\emergencystretch}{3em}
\newcommand{\N}{\mathbb{N}}
\newcommand{\Z}{\mathbb{Z}}
\newcommand{\Q}{\mathbb{Q}}
\newcommand{\R}{\mathbb{R}}
\newcommand{\C}{\mathbb{C}}
\newcommand{\F}{\mathbb{F}}
\newcommand{\A}{\mathbb{A}}
\newcommand{\PP}{\mathbb{P}}
\newcommand{\E}{\mathbb{E}}
\newcommand{\eps}{\varepsilon}
\newcommand{\dd}{\,\mathrm{d}}
\begin{document}
\section*{UnsolvedMath ID: 30006061}
\textbf{Source problem code:} catalogue entry
\subsection{Definitions and mathematical statement}
This is a statement-only reconstruction of the three conjectures recorded for UnsolvedMath ID \(30006061\). The supplied record does not specify the class of noncommutative del Pezzo surfaces, the base field, or several necessary conventions. The statements below retain those unresolved qualifications.

For a surface \(S\), write \(\mathcal D(S)\) for its intended derived category, whose precise definition must be supplied. In a \(k\)-linear triangulated category, an object \(E\) is exceptional if
\[
\operatorname{Hom}(E,E[i])\cong
\begin{cases}
k,&i=0,\\
0,&i\ne 0.
\end{cases}
\]
Here \(k\) denotes the base field when such a framework is specified; no field or characteristic is selected by the record. An exceptional collection is a finite ordered sequence \((E_1,\ldots,E_r)\) of exceptional objects such that
\[
\operatorname{Hom}(E_j,E_i[m])=0
\quad\text{for every }j>i\text{ and every }m\in\Z.
\]
It is full if its objects generate \(\mathcal D(S)\) as a triangulated category closed under direct summands. Mutation-equivalence means connection by a finite sequence of the intended left and right mutations of adjacent exceptional objects; whether shifts or other operations are additionally allowed is not specified.

The Grothendieck group \(K_0(\mathcal D(S))\) is generated by object classes \([E]\), subject to \([B]=[A]+[C]\) for each distinguished triangle \(A\to B\to C\to A[1]\). The record calls for the Grothendieck group, without specifying a numerical quotient.

\begin{enumerate}
\item \textbf{Exceptional objects and collections.} The source proposes that an exceptional object on a noncommutative del Pezzo surface is determined, up to the intended notion of isomorphism, by its Grothendieck-group class. Its proposed implication is
\[
[E]=[F]\quad\Longrightarrow\quad E\cong F
\]
for exceptional objects in the intended class. It further states that a marking is consequently equivalent to its image in the Grothendieck group, and that, for every such surface \(S\), any two full exceptional collections in \(\mathcal D(S)\) are related by mutations.

The eligible class of objects requires clarification: unrestricted derived objects satisfy \([E[2]]=[E]\), so the displayed implication needs a convention concerning shifts or a restriction on objects. No such convention is supplied. The definition of a marking and its Grothendieck-group image is also absent.

\item \textbf{Deformation of pure helices.} For every pure helix \(H\) on a commutative del Pezzo surface \(X\), and every noncommutative del Pezzo surface \(S\) in the intended deformation setting, the source proposes that \(H\) deforms to a pure helix on \(S\).

The record does not define a pure helix, its period, the relevant purity condition, the deformation family, or the compatibility required between \(X\) and \(S\). Thus the phrase ``in the intended deformation setting'' is an unresolved qualification, not an additional established hypothesis or a specified deformation construction.

\item \textbf{Classification by affine Weyl orbits.} Let \(\mathcal M\) denote the intended moduli stack of marked noncommutative del Pezzo surfaces, and let \(U(x)\) denote the underlying noncommutative surface of a geometric point \(x\) of \(\mathcal M\). Using the source notation, set
\[
W_{\mathrm{aff}}(\omega^\perp)
:=\omega^\perp\rtimes W(\omega^\perp).
\]
For geometric points \(x,y\) in the intended common moduli setting, the proposed assertion is
\[
U(x)\cong U(y)
\quad\Longleftrightarrow\quad
\exists g\in W_{\mathrm{aff}}(\omega^\perp)\text{ such that }y\simeq g\cdot x.
\]
Here \(\simeq\) means isomorphism of points in the moduli stack. The surfaces on the left are the underlying surfaces after forgetting the markings, as in the original passage. The ambient lattice, its pairing, the element \(\omega\), the associated Weyl group, the marking data, the group action, and the base for comparing geometric points are not defined in the record. The displayed formula therefore records the intended orbit criterion conditionally on those definitions and on the existence of the action.
\end{enumerate}

The supplied passage additionally reports that Conjecture 1 is known for commutative del Pezzo surfaces by Kuleshov--Orlov. It describes transporting a full exceptional collection from a commutative surface to a corresponding noncommutative surface by expressing it through mutations of a standard collection. It also reports an action of \(W_{\mathrm{aff}}(\omega^\perp)\) on noncommutative markings of a commutative del Pezzo surface, with the subgroup \(\omega^\perp\) acting by tensoring with line bundles, and says that Conjecture 2 would imply an affine Weyl group action on \(\mathcal M\). These are supplied contextual claims; the standard collection, correspondence of surfaces, and constructions of these actions remain unspecified here.
\subsection{Short English statement}
The source proposes three conjectures: exceptional objects are determined by their Grothendieck classes and full exceptional collections are connected by mutations; every pure helix deforms to the intended noncommutative surface; and underlying surfaces are isomorphic exactly when their marked moduli points belong to one affine Weyl orbit. Definitions of the surfaces, markings, helices, deformation setting, and group action require clarification, as does the treatment of shifts of exceptional objects. The supplied open status is dated curation, not an independently verified conclusion.
\subsection{Sources}
\begin{itemize}
\item Local UnsolvedMath registry, supplied complete source record for permanent ID \(30006061\), title ``Exceptional Collections on Noncommutative Del Pezzo Surfaces,'' exact source problem code \texttt{catalogue entry}. The record belongs to the Oberwolfach Reports Open Problems collection and gives proposed year 2024; no individual proposer is named.
\item Supplied report citation: \emph{MFO-RIMS Tandem Workshop: Algebraic Geometry and Noncommutative Projective Varieties (2025)}, Oberwolfach Reports, pp.\ 2551--2606, DOI \texttt{10.4171/owr/2024/44}. \url{https://doi.org/10.4171/owr/2024/44}. The literature-source label is ``Oberwolfach Report 44/2024.'' The supplied title year 2025 and report designation 44/2024 are retained without resolving their discrepancy.
\item Supplied literature source: Nordskova--Van den Bergh, \emph{NCCRs of cones over del Pezzo surfaces}. \url{https://arxiv.org/abs/2604.11319}.
\item Literature assessment dated exactly \texttt{2026-08-22}: the supplied curation labels the problem open and states that later mutation results for del Pezzo cones do not establish all three conjectures. This is supplied dated curation, not an independently verified literature assessment.
\item The record attributes the commutative case of Conjecture 1 to Kuleshov--Orlov, but supplies no separate bibliographic citation for that attribution. It also labels the statement \texttt{corrected\_verified} and says it was reconstructed and checked against the official report. These are source-record assertions; no source was independently consulted for this draft.
\end{itemize}
\end{document}
